\section{晚清的学校、家庭、环境与地方制度}
灾荒、铁路、学校、军队、公司、教案和新政共同改变了晚清地方社会的材料条件。本节以教育、劳动、家庭、环境与城市公共生活为问题，不把报刊舆论、改革法令或海关统计误认为全国经验；每个节点均保留其文书机构和社会覆盖范围。
\subsection{1870—1879：地方材料所能看见的尺度}
以下事件按账本时间排列；它们不是同质的社会样本，而是提示应到何种地点、机构和材料中继续追问的入口。
\hypertarget{social:EQ-1875-GUANGXU-ACCESSION}{}\paragraph{光绪元年（1875年）}光绪帝即位，慈禧、慈安继续掌握关键政治影响。从地方治理、身份、家庭与物质生活的角度，这一节点要求追问它在何处改变了资源、道路、组织或日常风险。核心来源为《清德宗实录》光绪元年相关条；宫中朱批奏折、内阁大库题本与《清史稿》相关志传。它首先留下的是官府分类、任命、处置或资源调拨的痕迹；要判断基层是否照办，仍须寻找县档、账簿、契约或后续案件。账本所示的研究方向是宫廷政治、制度运作与中央—地方信息链研究。可核的条件包括幼主继承使摄政与名分安排再度出现。应分层观察的后果是清末改革和对外危机将在这一权力格局下发展。证据界限：以光绪元年（1875年）的同期文书为定位起点；官修记录、奏报或外交文本服务特定机构立场，具体日期、规模、地方执行与对手视角须以独立材料复核。
\hypertarget{social:EQ-1875-YUNNAN-MUSLIM-WAR}{}\paragraph{光绪元年（1875年）}杜文秀政权覆亡前后，云南回民战争进入惨烈收束阶段。从地方治理、身份、家庭与物质生活的角度，这一节点要求追问它在何处改变了资源、道路、组织或日常风险。核心来源为《清德宗实录》光绪元年相关条；军机处档、战事方略及地方奏报。这类编年材料可以标定朝廷获知和叙述事件的时间，却往往压低妇女、儿童、佃户、流民和非文书精英的行动。账本所示的研究方向是战争动员、军需、战场暴力与地方社会研究。可核的条件包括地方军政、矿业贸易与族群冲突长期累积。应分层观察的后果是死亡、迁徙与复垦规模的估计须保持不确定性。证据界限：以光绪元年（1875年）的同期文书为定位起点；官修记录、奏报或外交文本服务特定机构立场，具体日期、规模、地方执行与对手视角须以独立材料复核。
\hypertarget{social:EQ-1876-CHEFOO-CONVENTION}{}\paragraph{光绪二年（1876年）}《烟台条约》签订，扩大英国在华通商和内地活动权利。从地方治理、身份、家庭与物质生活的角度，这一节点要求追问它在何处改变了资源、道路、组织或日常风险。核心来源为《清德宗实录》光绪二年相关条；《筹办夷务始末》相关年卷与中外照会、条约文本。它首先留下的是官府分类、任命、处置或资源调拨的痕迹；要判断基层是否照办，仍须寻找县档、账簿、契约或后续案件。账本所示的研究方向是条约文本、外交档案、口岸社会与国际法实践研究。可核的条件包括马嘉理事件与英国外交压力推动谈判。应分层观察的后果是条约权利的地区落实和地方反应需要逐案追踪。证据界限：以光绪二年（1876年）的同期文书为定位起点；官修记录、奏报或外交文本服务特定机构立场，具体日期、规模、地方执行与对手视角须以独立材料复核。
\hypertarget{social:EQ-1876-NORTH-CHINA-FAMINE}{}\paragraph{光绪二年（1876年）}丁戊奇荒在华北及西北造成广泛饥荒、迁徙和救济行动。从地方治理、身份、家庭与物质生活的角度，这一节点要求追问它在何处改变了资源、道路、组织或日常风险。核心来源为《清德宗实录》光绪二年相关条；地方志、督抚奏折及地方档案。这类编年材料可以标定朝廷获知和叙述事件的时间，却往往压低妇女、儿童、佃户、流民和非文书精英的行动。账本所示的研究方向是地方档案、迁徙、宗教、族群与社会动员研究。可核的条件包括旱灾、粮价、交通和行政能力共同塑造灾情。应分层观察的后果是伤亡总数争议极大，应按地区、材料与估算口径说明。证据界限：以光绪二年（1876年）的同期文书为定位起点；官修记录、奏报或外交文本服务特定机构立场，具体日期、规模、地方执行与对手视角须以独立材料复核。
\hypertarget{social:EQ-1877-FAMINE-RELIEF-NETWORKS}{}\paragraph{光绪三年（1877年）}赈济网络在灾区展开，官办、绅办和传教团体的救助并存。从地方治理、身份、家庭与物质生活的角度，这一节点要求追问它在何处改变了资源、道路、组织或日常风险。核心来源为《清德宗实录》光绪三年相关条；地方志、督抚奏折及地方档案。这类编年材料可以标定朝廷获知和叙述事件的时间，却往往压低妇女、儿童、佃户、流民和非文书精英的行动。账本所示的研究方向是地方档案、迁徙、宗教、族群与社会动员研究。可核的条件包括灾害规模超出常规仓储和地方财政能力。应分层观察的后果是救助记录反映组织覆盖，不等于所有灾民获救。证据界限：以光绪三年（1877年）的同期文书为定位起点；官修记录、奏报或外交文本服务特定机构立场，具体日期、规模、地方执行与对手视角须以独立材料复核。
\hypertarget{social:EQ-1878-Zuo-ZONGTANG-XINJIANG}{}\paragraph{光绪四年（1878年）}左宗棠军收复新疆大部，阿古柏政权崩解。从地方治理、身份、家庭与物质生活的角度，这一节点要求追问它在何处改变了资源、道路、组织或日常风险。核心来源为《清德宗实录》光绪四年相关条；军机处档、战事方略及地方奏报。这类编年材料可以标定朝廷获知和叙述事件的时间，却往往压低妇女、儿童、佃户、流民和非文书精英的行动。账本所示的研究方向是战争动员、军需、战场暴力与地方社会研究。可核的条件包括清廷集中湘军、饷源和西北运输资源进行反攻。应分层观察的后果是军事收复与后续行政、绿洲社会重建必须分开叙述。证据界限：以光绪四年（1878年）的同期文书为定位起点；官修记录、奏报或外交文本服务特定机构立场，具体日期、规模、地方执行与对手视角须以独立材料复核。
\hypertarget{social:EQ-1878-XINJIANG-PROVINCE-DEBATE}{}\paragraph{光绪四年（1878年）}朝廷讨论新疆设省、军政与财政安排。从地方治理、身份、家庭与物质生活的角度，这一节点要求追问它在何处改变了资源、道路、组织或日常风险。核心来源为《清德宗实录》光绪四年相关条；宫中朱批奏折、内阁大库题本与《清史稿》相关志传。它首先留下的是官府分类、任命、处置或资源调拨的痕迹；要判断基层是否照办，仍须寻找县档、账簿、契约或后续案件。账本所示的研究方向是宫廷政治、制度运作与中央—地方信息链研究。可核的条件包括战后长期驻防成本和俄国边境压力要求重建制度。应分层观察的后果是设省方案反映中央目标，地方实施尚需持续资源。证据界限：以光绪四年（1878年）的同期文书为定位起点；官修记录、奏报或外交文本服务特定机构立场，具体日期、规模、地方执行与对手视角须以独立材料复核。
\hypertarget{social:EQ-1879-RYUKYU-ANNEXATION}{}\paragraph{光绪五年（1879年）}日本吞并琉球，清廷的交涉未能恢复旧有册封关系。从地方治理、身份、家庭与物质生活的角度，这一节点要求追问它在何处改变了资源、道路、组织或日常风险。核心来源为《清德宗实录》光绪五年相关条；《筹办夷务始末》相关年卷与中外照会、条约文本。它首先留下的是官府分类、任命、处置或资源调拨的痕迹；要判断基层是否照办，仍须寻找县档、账簿、契约或后续案件。账本所示的研究方向是条约文本、外交档案、口岸社会与国际法实践研究。可核的条件包括日本国家建构与东亚秩序转变相互作用。应分层观察的后果是清方外交文书与琉球、日本材料对权利表述不同。证据界限：以光绪五年（1879年）的同期文书为定位起点；官修记录、奏报或外交文本服务特定机构立场，具体日期、规模、地方执行与对手视角须以独立材料复核。
\subsection{1880—1889：地方材料所能看见的尺度}
以下事件按账本时间排列；它们不是同质的社会样本，而是提示应到何种地点、机构和材料中继续追问的入口。
\hypertarget{social:EQ-1880-ILI-NEGOTIATIONS}{}\paragraph{光绪六年（1880年）}清俄围绕伊犁归还和边界条件持续谈判。从地方治理、身份、家庭与物质生活的角度，这一节点要求追问它在何处改变了资源、道路、组织或日常风险。核心来源为《清德宗实录》光绪六年相关条；《筹办夷务始末》相关年卷与中外照会、条约文本。它首先留下的是官府分类、任命、处置或资源调拨的痕迹；要判断基层是否照办，仍须寻找县档、账簿、契约或后续案件。账本所示的研究方向是条约文本、外交档案、口岸社会与国际法实践研究。可核的条件包括俄军占据伊犁后，清廷在军事与财政限制下寻求外交解决。应分层观察的后果是边界线、赔款和贸易条款须逐项区分。证据界限：以光绪六年（1880年）的同期文书为定位起点；官修记录、奏报或外交文本服务特定机构立场，具体日期、规模、地方执行与对手视角须以独立材料复核。
\hypertarget{social:EQ-1881-SAINT-PETERSBURG-TREATY}{}\paragraph{光绪七年（1881年）}《中俄伊犁条约》签订，伊犁大部归还清廷并附带领土、赔款和贸易条件。从地方治理、身份、家庭与物质生活的角度，这一节点要求追问它在何处改变了资源、道路、组织或日常风险。核心来源为《清德宗实录》光绪七年相关条；《筹办夷务始末》相关年卷与中外照会、条约文本。它首先留下的是官府分类、任命、处置或资源调拨的痕迹；要判断基层是否照办，仍须寻找县档、账簿、契约或后续案件。账本所示的研究方向是条约文本、外交档案、口岸社会与国际法实践研究。可核的条件包括长期谈判与西北军事实力恢复共同作用。应分层观察的后果是条约地图与实际边界管理不是同一问题。证据界限：以光绪七年（1881年）的同期文书为定位起点；官修记录、奏报或外交文本服务特定机构立场，具体日期、规模、地方执行与对手视角须以独立材料复核。
\hypertarget{social:EQ-1882-IMO-INCIDENT}{}\paragraph{光绪八年（1882年）}朝鲜壬午兵变后清军入朝，日本也扩大影响。从地方治理、身份、家庭与物质生活的角度，这一节点要求追问它在何处改变了资源、道路、组织或日常风险。核心来源为《清德宗实录》光绪八年相关条；《筹办夷务始末》相关年卷与中外照会、条约文本。它首先留下的是官府分类、任命、处置或资源调拨的痕迹；要判断基层是否照办，仍须寻找县档、账簿、契约或后续案件。账本所示的研究方向是条约文本、外交档案、口岸社会与国际法实践研究。可核的条件包括朝鲜内政、军制改革与列强竞争交织。应分层观察的后果是清、朝、日文献中的“保护”与控制含义不同。证据界限：以光绪八年（1882年）的同期文书为定位起点；官修记录、奏报或外交文本服务特定机构立场，具体日期、规模、地方执行与对手视角须以独立材料复核。
\hypertarget{social:EQ-1883-SINO-FRENCH-WAR-BEGIN}{}\paragraph{光绪九年（1883年）}中法因越南北部和宗藩关系冲突升级为战争。从地方治理、身份、家庭与物质生活的角度，这一节点要求追问它在何处改变了资源、道路、组织或日常风险。核心来源为《清德宗实录》光绪九年相关条；军机处档、战事方略及地方奏报。这类编年材料可以标定朝廷获知和叙述事件的时间，却往往压低妇女、儿童、佃户、流民和非文书精英的行动。账本所示的研究方向是战争动员、军需、战场暴力与地方社会研究。可核的条件包括法国殖民扩张与清廷在越南的传统关系碰撞。应分层观察的后果是战事跨越越南、台湾和华南，不能只从北京外交看。证据界限：以光绪九年（1883年）的同期文书为定位起点；官修记录、奏报或外交文本服务特定机构立场，具体日期、规模、地方执行与对手视角须以独立材料复核。
\hypertarget{social:EQ-1884-FUZHOU-NAVAL-BATTLE}{}\paragraph{光绪十年（1884年）}马江海战中福建水师遭法国舰队重创。从地方治理、身份、家庭与物质生活的角度，这一节点要求追问它在何处改变了资源、道路、组织或日常风险。核心来源为《清德宗实录》光绪十年相关条；军机处档、战事方略及地方奏报。这类编年材料可以标定朝廷获知和叙述事件的时间，却往往压低妇女、儿童、佃户、流民和非文书精英的行动。账本所示的研究方向是战争动员、军需、战场暴力与地方社会研究。可核的条件包括清军舰队训练、指挥与武器体系面临现代海战挑战。应分层观察的后果是伤亡和舰只损失应依据中法记录和船政资料核验。证据界限：以光绪十年（1884年）的同期文书为定位起点；官修记录、奏报或外交文本服务特定机构立场，具体日期、规模、地方执行与对手视角须以独立材料复核。
\hypertarget{social:EQ-1884-TAIWAN-BLOCKADE}{}\paragraph{光绪十年（1884年）}法军进攻台湾北部并封锁，台湾成为中法战争重要战场。从地方治理、身份、家庭与物质生活的角度，这一节点要求追问它在何处改变了资源、道路、组织或日常风险。核心来源为《清德宗实录》光绪十年相关条；军机处档、战事方略及地方奏报。这类编年材料可以标定朝廷获知和叙述事件的时间，却往往压低妇女、儿童、佃户、流民和非文书精英的行动。账本所示的研究方向是战争动员、军需、战场暴力与地方社会研究。可核的条件包括台湾港口、煤矿和海峡航路具有战略价值。应分层观察的后果是地方抵抗、商业中断与原住民地区情况需分区考察。证据界限：以光绪十年（1884年）的同期文书为定位起点；官修记录、奏报或外交文本服务特定机构立场，具体日期、规模、地方执行与对手视角须以独立材料复核。
\hypertarget{social:EQ-1885-TIANJIN-CHINA-FRANCE-TREATY}{}\paragraph{光绪十一年（1885年）}《中法新约》确认越南脱离清朝宗藩体系，战争结束。从地方治理、身份、家庭与物质生活的角度，这一节点要求追问它在何处改变了资源、道路、组织或日常风险。核心来源为《清德宗实录》光绪十一年相关条；《筹办夷务始末》相关年卷与中外照会、条约文本。它首先留下的是官府分类、任命、处置或资源调拨的痕迹；要判断基层是否照办，仍须寻找县档、账簿、契约或后续案件。账本所示的研究方向是条约文本、外交档案、口岸社会与国际法实践研究。可核的条件包括军事与财政压力推动双方议和。应分层观察的后果是条约确认与边界勘界、华南地方影响应分别处理。证据界限：以光绪十一年（1885年）的同期文书为定位起点；官修记录、奏报或外交文本服务特定机构立场，具体日期、规模、地方执行与对手视角须以独立材料复核。
\hypertarget{social:EQ-1885-TAIWAN-PROVINCE}{}\paragraph{光绪十一年（1885年）}清廷决定将台湾建省，加强行政、军防和基础设施经营。从地方治理、身份、家庭与物质生活的角度，这一节点要求追问它在何处改变了资源、道路、组织或日常风险。核心来源为《清德宗实录》光绪十一年相关条；宫中朱批奏折、内阁大库题本与《清史稿》相关志传。它首先留下的是官府分类、任命、处置或资源调拨的痕迹；要判断基层是否照办，仍须寻找县档、账簿、契约或后续案件。账本所示的研究方向是宫廷政治、制度运作与中央—地方信息链研究。可核的条件包括中法战争显示台湾海防与行政地位的重要性。应分层观察的后果是建省不等于全岛被均质治理，东部和原住民地区尤需谨慎。证据界限：以光绪十一年（1885年）的同期文书为定位起点；官修记录、奏报或外交文本服务特定机构立场，具体日期、规模、地方执行与对手视角须以独立材料复核。
\hypertarget{social:EQ-1886-BEIYANG-NAVY-EXPANSION}{}\paragraph{光绪十二年（1886年）}北洋海军和沿海防务建设持续推进。从地方治理、身份、家庭与物质生活的角度，这一节点要求追问它在何处改变了资源、道路、组织或日常风险。核心来源为《清德宗实录》光绪十二年相关条；军机处档、战事方略及地方奏报。这类编年材料可以标定朝廷获知和叙述事件的时间，却往往压低妇女、儿童、佃户、流民和非文书精英的行动。账本所示的研究方向是战争动员、军需、战场暴力与地方社会研究。可核的条件包括中法战争和日本扩张促使海军现代化。应分层观察的后果是购舰数量不等于训练、维护、弹药和指挥能力。证据界限：以光绪十二年（1886年）的同期文书为定位起点；官修记录、奏报或外交文本服务特定机构立场，具体日期、规模、地方执行与对手视角须以独立材料复核。
\hypertarget{social:EQ-1887-TAIWAN-RAILWAY-PLANNING}{}\paragraph{光绪十三年（1887年）}台湾巡抚刘铭传推动铁路、电报、矿务和新式行政项目。从地方治理、身份、家庭与物质生活的角度，这一节点要求追问它在何处改变了资源、道路、组织或日常风险。核心来源为《清德宗实录》光绪十三年相关条；户部奏销、盐政、河工或海关档案。它首先留下的是官府分类、任命、处置或资源调拨的痕迹；要判断基层是否照办，仍须寻找县档、账簿、契约或后续案件。账本所示的研究方向是财政账册、市场网络、灾害环境与区域经济研究。可核的条件包括建省后需要连接港口、城市和军防节点。应分层观察的后果是项目经费、路段使用与社会影响需具体核验。证据界限：以光绪十三年（1887年）的同期文书为定位起点；官修记录、奏报或外交文本服务特定机构立场，具体日期、规模、地方执行与对手视角须以独立材料复核。
\hypertarget{social:EQ-1888-NAVAL-AND-ARMY-REFORM}{}\paragraph{光绪十四年（1888年）}清廷讨论海军、陆军训练和军工体系的经费与制度问题。从地方治理、身份、家庭与物质生活的角度，这一节点要求追问它在何处改变了资源、道路、组织或日常风险。核心来源为《清德宗实录》光绪十四年相关条；宫中朱批奏折、内阁大库题本与《清史稿》相关志传。它首先留下的是官府分类、任命、处置或资源调拨的痕迹；要判断基层是否照办，仍须寻找县档、账簿、契约或后续案件。账本所示的研究方向是宫廷政治、制度运作与中央—地方信息链研究。可核的条件包括列强战争压力暴露既有军制不足。应分层观察的后果是改革奏议与实际编制、操练之间常有距离。证据界限：以光绪十四年（1888年）的同期文书为定位起点；官修记录、奏报或外交文本服务特定机构立场，具体日期、规模、地方执行与对手视角须以独立材料复核。
\hypertarget{social:EQ-1889-GUANGXU-MARRIAGE}{}\paragraph{光绪十五年（1889年）}光绪帝大婚并开始形式上的亲政安排。从地方治理、身份、家庭与物质生活的角度，这一节点要求追问它在何处改变了资源、道路、组织或日常风险。核心来源为《清德宗实录》光绪十五年相关条；宫中朱批奏折、内阁大库题本与《清史稿》相关志传。它首先留下的是官府分类、任命、处置或资源调拨的痕迹；要判断基层是否照办，仍须寻找县档、账簿、契约或后续案件。账本所示的研究方向是宫廷政治、制度运作与中央—地方信息链研究。可核的条件包括皇帝成年要求完成宫廷礼制程序。应分层观察的后果是慈禧仍对重要决策保有深刻影响。证据界限：以光绪十五年（1889年）的同期文书为定位起点；官修记录、奏报或外交文本服务特定机构立场，具体日期、规模、地方执行与对手视角须以独立材料复核。
\subsection{1890—1899：地方材料所能看见的尺度}
以下事件按账本时间排列；它们不是同质的社会样本，而是提示应到何种地点、机构和材料中继续追问的入口。
\hypertarget{social:EQ-1890-KOREAN-POLITICAL-CONFLICT}{}\paragraph{光绪十六年（1890年）}朝鲜内政纷争持续牵动清日两国在半岛的影响力竞争。从地方治理、身份、家庭与物质生活的角度，这一节点要求追问它在何处改变了资源、道路、组织或日常风险。核心来源为《清德宗实录》光绪十六年相关条；《筹办夷务始末》相关年卷与中外照会、条约文本。它首先留下的是官府分类、任命、处置或资源调拨的痕迹；要判断基层是否照办，仍须寻找县档、账簿、契约或后续案件。账本所示的研究方向是条约文本、外交档案、口岸社会与国际法实践研究。可核的条件包括壬午事变后清日均扩大驻军与交涉。应分层观察的后果是朝鲜行动者不是外部竞争的被动对象，应纳入叙述。证据界限：以光绪十六年（1890年）的同期文书为定位起点；官修记录、奏报或外交文本服务特定机构立场，具体日期、规模、地方执行与对手视角须以独立材料复核。
\hypertarget{social:EQ-1891-MISSIONARY-CONFLICTS}{}\paragraph{光绪十七年（1891年）}各地教案和宗教纠纷继续引发地方治理与外交压力。从地方治理、身份、家庭与物质生活的角度，这一节点要求追问它在何处改变了资源、道路、组织或日常风险。核心来源为《清德宗实录》光绪十七年相关条；地方志、督抚奏折及地方档案。这类编年材料可以标定朝廷获知和叙述事件的时间，却往往压低妇女、儿童、佃户、流民和非文书精英的行动。账本所示的研究方向是地方档案、迁徙、宗教、族群与社会动员研究。可核的条件包括土地、诉讼、保护权和谣言与宗教身份交织。应分层观察的后果是“教案”包含不同类型纠纷，不能单一化。证据界限：以光绪十七年（1891年）的同期文书为定位起点；官修记录、奏报或外交文本服务特定机构立场，具体日期、规模、地方执行与对手视角须以独立材料复核。
\hypertarget{social:EQ-1892-RAILWAY-DEBATE}{}\paragraph{光绪十八年（1892年）}官员、商人和士绅围绕铁路修筑、主权与资金展开争论。从地方治理、身份、家庭与物质生活的角度，这一节点要求追问它在何处改变了资源、道路、组织或日常风险。核心来源为《清德宗实录》光绪十八年相关条；户部奏销、盐政、河工或海关档案。它首先留下的是官府分类、任命、处置或资源调拨的痕迹；要判断基层是否照办，仍须寻找县档、账簿、契约或后续案件。账本所示的研究方向是财政账册、市场网络、灾害环境与区域经济研究。可核的条件包括交通技术、外债和地方利益引发新问题。应分层观察的后果是赞成或反对文本不等于实际线路建设和公众意见。证据界限：以光绪十八年（1892年）的同期文书为定位起点；官修记录、奏报或外交文本服务特定机构立场，具体日期、规模、地方执行与对手视角须以独立材料复核。
\hypertarget{social:EQ-1893-KOREAN-REFORM-CRISIS}{}\paragraph{光绪十九年（1893年）}朝鲜改革与政治危机使清日关系持续紧张。从地方治理、身份、家庭与物质生活的角度，这一节点要求追问它在何处改变了资源、道路、组织或日常风险。核心来源为《清德宗实录》光绪十九年相关条；《筹办夷务始末》相关年卷与中外照会、条约文本。它首先留下的是官府分类、任命、处置或资源调拨的痕迹；要判断基层是否照办，仍须寻找县档、账簿、契约或后续案件。账本所示的研究方向是条约文本、外交档案、口岸社会与国际法实践研究。可核的条件包括半岛内政和驻军安排缺乏稳定共识。应分层观察的后果是清日文书各自服务战略目的，需与朝鲜材料互读。证据界限：以光绪十九年（1893年）的同期文书为定位起点；官修记录、奏报或外交文本服务特定机构立场，具体日期、规模、地方执行与对手视角须以独立材料复核。
\hypertarget{social:EQ-1894-DONGHAK-PEASANT-WAR}{}\paragraph{光绪二十年（1894年）}东学农民战争促使清日两国派兵入朝，危机迅速国际化。从地方治理、身份、家庭与物质生活的角度，这一节点要求追问它在何处改变了资源、道路、组织或日常风险。核心来源为《清德宗实录》光绪二十年相关条；军机处档、战事方略及地方奏报。这类编年材料可以标定朝廷获知和叙述事件的时间，却往往压低妇女、儿童、佃户、流民和非文书精英的行动。账本所示的研究方向是战争动员、军需、战场暴力与地方社会研究。可核的条件包括朝鲜社会改革诉求与政府镇压失败交织。应分层观察的后果是农民运动不应只作为清日战争的前奏叙述。证据界限：以光绪二十年（1894年）的同期文书为定位起点；官修记录、奏报或外交文本服务特定机构立场，具体日期、规模、地方执行与对手视角须以独立材料复核。
\hypertarget{social:EQ-1894-FIRST-SINO-JAPANESE-WAR}{}\paragraph{光绪二十年（1894年）}甲午战争爆发，清日围绕朝鲜和区域秩序展开海陆战。从地方治理、身份、家庭与物质生活的角度，这一节点要求追问它在何处改变了资源、道路、组织或日常风险。核心来源为《清德宗实录》光绪二十年相关条；军机处档、战事方略及地方奏报。这类编年材料可以标定朝廷获知和叙述事件的时间，却往往压低妇女、儿童、佃户、流民和非文书精英的行动。账本所示的研究方向是战争动员、军需、战场暴力与地方社会研究。可核的条件包括驻军竞争与外交谈判破裂升级为战争。应分层观察的后果是战场损失、指挥责任与民众经验须多源核验。证据界限：以光绪二十年（1894年）的同期文书为定位起点；官修记录、奏报或外交文本服务特定机构立场，具体日期、规模、地方执行与对手视角须以独立材料复核。
\hypertarget{social:EQ-1894-YALU-NAVAL-BATTLE}{}\paragraph{光绪二十年（1894年）}黄海海战后北洋舰队受创，海军力量对比进一步恶化。从地方治理、身份、家庭与物质生活的角度，这一节点要求追问它在何处改变了资源、道路、组织或日常风险。核心来源为《清德宗实录》光绪二十年相关条；军机处档、战事方略及地方奏报。这类编年材料可以标定朝廷获知和叙述事件的时间，却往往压低妇女、儿童、佃户、流民和非文书精英的行动。账本所示的研究方向是战争动员、军需、战场暴力与地方社会研究。可核的条件包括舰艇性能、训练、战术和指挥因素共同作用。应分层观察的后果是单一“装备落后”解释无法涵盖战役复杂性。证据界限：以光绪二十年（1894年）的同期文书为定位起点；官修记录、奏报或外交文本服务特定机构立场，具体日期、规模、地方执行与对手视角须以独立材料复核。
\hypertarget{social:EQ-1895-SHIMONOSEKI-TREATY}{}\paragraph{光绪二十一年（1895年）}《马关条约》签订，清割让台湾、澎湖并承认朝鲜独立，赔款与通商条件加重。从地方治理、身份、家庭与物质生活的角度，这一节点要求追问它在何处改变了资源、道路、组织或日常风险。核心来源为《清德宗实录》光绪二十一年相关条；《筹办夷务始末》相关年卷与中外照会、条约文本。它首先留下的是官府分类、任命、处置或资源调拨的痕迹；要判断基层是否照办，仍须寻找县档、账簿、契约或后续案件。账本所示的研究方向是条约文本、外交档案、口岸社会与国际法实践研究。可核的条件包括清军战败和日本军事压力迫使议和。应分层观察的后果是条约的签署、割让交接和地方抵抗是不同过程。证据界限：以光绪二十一年（1895年）的同期文书为定位起点；官修记录、奏报或外交文本服务特定机构立场，具体日期、规模、地方执行与对手视角须以独立材料复核。
\hypertarget{social:EQ-1895-TAIWAN-REPUBLIC-RESISTANCE}{}\paragraph{光绪二十一年（1895年）}台湾割让后地方力量建立台湾民主国并抵抗日军登陆。从地方治理、身份、家庭与物质生活的角度，这一节点要求追问它在何处改变了资源、道路、组织或日常风险。核心来源为《清德宗实录》光绪二十一年相关条；军机处档、战事方略及地方奏报。这类编年材料可以标定朝廷获知和叙述事件的时间，却往往压低妇女、儿童、佃户、流民和非文书精英的行动。账本所示的研究方向是战争动员、军需、战场暴力与地方社会研究。可核的条件包括条约割让与地方社会的政治、军事反应不一致。应分层观察的后果是抵抗组织与普通居民处境不能由单一民族叙事概括。证据界限：以光绪二十一年（1895年）的同期文书为定位起点；官修记录、奏报或外交文本服务特定机构立场，具体日期、规模、地方执行与对手视角须以独立材料复核。
\hypertarget{social:EQ-1895-TRIPLE-INTERVENTION}{}\paragraph{光绪二十一年（1895年）}三国干涉还辽改变战后列强竞争和清廷外交处境。从地方治理、身份、家庭与物质生活的角度，这一节点要求追问它在何处改变了资源、道路、组织或日常风险。核心来源为《清德宗实录》光绪二十一年相关条；《筹办夷务始末》相关年卷与中外照会、条约文本。它首先留下的是官府分类、任命、处置或资源调拨的痕迹；要判断基层是否照办，仍须寻找县档、账簿、契约或后续案件。账本所示的研究方向是条约文本、外交档案、口岸社会与国际法实践研究。可核的条件包括列强在东北和东亚的利益冲突加剧。应分层观察的后果是干涉并非出于中立调停，需置于帝国竞争中理解。证据界限：以光绪二十一年（1895年）的同期文书为定位起点；官修记录、奏报或外交文本服务特定机构立场，具体日期、规模、地方执行与对手视角须以独立材料复核。
\hypertarget{social:EQ-1896-LI-LOBANOV-TREATY}{}\paragraph{光绪二十二年（1896年）}《中俄密约》及东清铁路安排加强俄国在东北的影响。从地方治理、身份、家庭与物质生活的角度，这一节点要求追问它在何处改变了资源、道路、组织或日常风险。核心来源为《清德宗实录》光绪二十二年相关条；《筹办夷务始末》相关年卷与中外照会、条约文本。它首先留下的是官府分类、任命、处置或资源调拨的痕迹；要判断基层是否照办，仍须寻找县档、账簿、契约或后续案件。账本所示的研究方向是条约文本、外交档案、口岸社会与国际法实践研究。可核的条件包括清廷寻求对日后的外交支撑，俄国谋求远东利益。应分层观察的后果是秘密条约的文本、执行与地方影响需分层说明。证据界限：以光绪二十二年（1896年）的同期文书为定位起点；官修记录、奏报或外交文本服务特定机构立场，具体日期、规模、地方执行与对手视角须以独立材料复核。
\hypertarget{social:EQ-1896-MODERN-EDUCATION-EXPANSION}{}\paragraph{光绪二十二年（1896年）}甲午战后新式学堂、译书与留学讨论加速。从地方治理、身份、家庭与物质生活的角度，这一节点要求追问它在何处改变了资源、道路、组织或日常风险。核心来源为《清德宗实录》光绪二十二年相关条；内阁档、文集或报刊相关记录。这类编年材料可以标定朝廷获知和叙述事件的时间，却往往压低妇女、儿童、佃户、流民和非文书精英的行动。账本所示的研究方向是知识生产、教育、宗教与公共传播研究。可核的条件包括战败促使士大夫和官员反思军政、教育与技术。应分层观察的后果是教育项目的地域覆盖和入学规模不宜夸大。证据界限：以光绪二十二年（1896年）的同期文书为定位起点；官修记录、奏报或外交文本服务特定机构立场，具体日期、规模、地方执行与对手视角须以独立材料复核。
\hypertarget{social:EQ-1897-JIAOZHOU-BAY-OCCUPATION}{}\paragraph{光绪二十三年（1897年）}德国以巨野教案为由占领胶州湾，清廷面临新的租借压力。从地方治理、身份、家庭与物质生活的角度，这一节点要求追问它在何处改变了资源、道路、组织或日常风险。核心来源为《清德宗实录》光绪二十三年相关条；《筹办夷务始末》相关年卷与中外照会、条约文本。它首先留下的是官府分类、任命、处置或资源调拨的痕迹；要判断基层是否照办，仍须寻找县档、账簿、契约或后续案件。账本所示的研究方向是条约文本、外交档案、口岸社会与国际法实践研究。可核的条件包括列强借教案和军舰扩张在华据点。应分层观察的后果是教案、军事占领和租借谈判应分别记录。证据界限：以光绪二十三年（1897年）的同期文书为定位起点；官修记录、奏报或外交文本服务特定机构立场，具体日期、规模、地方执行与对手视角须以独立材料复核。
\hypertarget{social:EQ-1898-LEASED-TERRITORIES}{}\paragraph{光绪二十四年（1898年）}俄、英、法、德相继取得租借地或势力范围安排。从地方治理、身份、家庭与物质生活的角度，这一节点要求追问它在何处改变了资源、道路、组织或日常风险。核心来源为《清德宗实录》光绪二十四年相关条；《筹办夷务始末》相关年卷与中外照会、条约文本。它首先留下的是官府分类、任命、处置或资源调拨的痕迹；要判断基层是否照办，仍须寻找县档、账簿、契约或后续案件。账本所示的研究方向是条约文本、外交档案、口岸社会与国际法实践研究。可核的条件包括甲午后清廷军事弱势与列强竞争叠加。应分层观察的后果是“势力范围”并不自动等于完整主权转移。证据界限：以光绪二十四年（1898年）的同期文书为定位起点；官修记录、奏报或外交文本服务特定机构立场，具体日期、规模、地方执行与对手视角须以独立材料复核。
\hypertarget{social:EQ-1898-HUNDRED-DAYS-REFORM}{}\paragraph{光绪二十四年（1898年）}光绪帝颁布变法诏令，推动官制、教育、经济与军事改革。从地方治理、身份、家庭与物质生活的角度，这一节点要求追问它在何处改变了资源、道路、组织或日常风险。核心来源为《清德宗实录》光绪二十四年相关条；宫中朱批奏折、内阁大库题本与《清史稿》相关志传。它首先留下的是官府分类、任命、处置或资源调拨的痕迹；要判断基层是否照办，仍须寻找县档、账簿、契约或后续案件。账本所示的研究方向是宫廷政治、制度运作与中央—地方信息链研究。可核的条件包括甲午战败和维新思潮使改革诉求集中爆发。应分层观察的后果是诏令与实际实施之间时间短、阻力大，须区别方案与结果。证据界限：以光绪二十四年（1898年）的同期文书为定位起点；官修记录、奏报或外交文本服务特定机构立场，具体日期、规模、地方执行与对手视角须以独立材料复核。
\hypertarget{social:EQ-1898-WUXU-COUP}{}\paragraph{光绪二十四年（1898年）}慈禧发动政变，光绪帝被幽禁，戊戌变法主要措施被撤销或改写。从地方治理、身份、家庭与物质生活的角度，这一节点要求追问它在何处改变了资源、道路、组织或日常风险。核心来源为《清德宗实录》光绪二十四年相关条；宫中朱批奏折、内阁大库题本与《清史稿》相关志传。它首先留下的是官府分类、任命、处置或资源调拨的痕迹；要判断基层是否照办，仍须寻找县档、账簿、契约或后续案件。账本所示的研究方向是宫廷政治、制度运作与中央—地方信息链研究。可核的条件包括改革触及官僚利益、皇权关系与军事权力。应分层观察的后果是政变叙述中的人物动机和密谋细节需谨慎处理。证据界限：以光绪二十四年（1898年）的同期文书为定位起点；官修记录、奏报或外交文本服务特定机构立场，具体日期、规模、地方执行与对手视角须以独立材料复核。
\hypertarget{social:EQ-1899-BOXER-UPRISING}{}\paragraph{光绪二十五年（1899年）}义和团在华北扩展，反教、反洋冲突加剧。从地方治理、身份、家庭与物质生活的角度，这一节点要求追问它在何处改变了资源、道路、组织或日常风险。核心来源为《清德宗实录》光绪二十五年相关条；地方志、督抚奏折及地方档案。这类编年材料可以标定朝廷获知和叙述事件的时间，却往往压低妇女、儿童、佃户、流民和非文书精英的行动。账本所示的研究方向是地方档案、迁徙、宗教、族群与社会动员研究。可核的条件包括灾荒、地方武术结社、教案和列强压力相互作用。应分层观察的后果是“义和团”内部成分复杂，不能以单一意识形态概括。证据界限：以光绪二十五年（1899年）的同期文书为定位起点；官修记录、奏报或外交文本服务特定机构立场，具体日期、规模、地方执行与对手视角须以独立材料复核。
\subsection{1900—1909：地方材料所能看见的尺度}
以下事件按账本时间排列；它们不是同质的社会样本，而是提示应到何种地点、机构和材料中继续追问的入口。
\hypertarget{social:EQ-1900-BOXER-WAR}{}\paragraph{光绪二十六年（1900年）}慈禧对列强宣战，八国联军入侵并占领北京。从地方治理、身份、家庭与物质生活的角度，这一节点要求追问它在何处改变了资源、道路、组织或日常风险。核心来源为《清德宗实录》光绪二十六年相关条；军机处档、战事方略及地方奏报。这类编年材料可以标定朝廷获知和叙述事件的时间，却往往压低妇女、儿童、佃户、流民和非文书精英的行动。账本所示的研究方向是战争动员、军需、战场暴力与地方社会研究。可核的条件包括朝廷内部政策摇摆、义和团动员与外交危机相互激化。应分层观察的后果是战争责任、暴力与劫掠需同时纳入中外和地方材料。证据界限：以光绪二十六年（1900年）的同期文书为定位起点；官修记录、奏报或外交文本服务特定机构立场，具体日期、规模、地方执行与对手视角须以独立材料复核。
\hypertarget{social:EQ-1900-SOUTHEAST-MUTUAL-PROTECTION}{}\paragraph{光绪二十六年（1900年）}东南互保使部分督抚与列强协商维持地方秩序。从地方治理、身份、家庭与物质生活的角度，这一节点要求追问它在何处改变了资源、道路、组织或日常风险。核心来源为《清德宗实录》光绪二十六年相关条；宫中朱批奏折、内阁大库题本与《清史稿》相关志传。它首先留下的是官府分类、任命、处置或资源调拨的痕迹；要判断基层是否照办，仍须寻找县档、账簿、契约或后续案件。账本所示的研究方向是宫廷政治、制度运作与中央—地方信息链研究。可核的条件包括中央战争命令与地方对商贸、治安的现实考量冲突。应分层观察的后果是“互保”并非整齐统一的地方独立，须逐省考察。证据界限：以光绪二十六年（1900年）的同期文书为定位起点；官修记录、奏报或外交文本服务特定机构立场，具体日期、规模、地方执行与对手视角须以独立材料复核。
\hypertarget{social:EQ-1901-BOXER-PROTOCOL}{}\paragraph{光绪二十七年（1901年）}《辛丑条约》签订，赔款、驻军和外交条件进一步限制清廷。从地方治理、身份、家庭与物质生活的角度，这一节点要求追问它在何处改变了资源、道路、组织或日常风险。核心来源为《清德宗实录》光绪二十七年相关条；《筹办夷务始末》相关年卷与中外照会、条约文本。它首先留下的是官府分类、任命、处置或资源调拨的痕迹；要判断基层是否照办，仍须寻找县档、账簿、契约或后续案件。账本所示的研究方向是条约文本、外交档案、口岸社会与国际法实践研究。可核的条件包括北京陷落和议和谈判使清廷处于弱势。应分层观察的后果是赔款总额、支付机制与地方税收后果应分开分析。证据界限：以光绪二十七年（1901年）的同期文书为定位起点；官修记录、奏报或外交文本服务特定机构立场，具体日期、规模、地方执行与对手视角须以独立材料复核。
\hypertarget{social:EQ-1901-ZONGLI-YAMEN-REPLACED}{}\paragraph{光绪二十七年（1901年）}总理衙门改为外务部，晚清外交机构进一步制度化。从地方治理、身份、家庭与物质生活的角度，这一节点要求追问它在何处改变了资源、道路、组织或日常风险。核心来源为《清德宗实录》光绪二十七年相关条；宫中朱批奏折、内阁大库题本与《清史稿》相关志传。它首先留下的是官府分类、任命、处置或资源调拨的痕迹；要判断基层是否照办，仍须寻找县档、账簿、契约或后续案件。账本所示的研究方向是宫廷政治、制度运作与中央—地方信息链研究。可核的条件包括战后改革要求重组对外事务处理。应分层观察的后果是机构改名不等于外交权力和信息网络立即现代化。证据界限：以光绪二十七年（1901年）的同期文书为定位起点；官修记录、奏报或外交文本服务特定机构立场，具体日期、规模、地方执行与对手视角须以独立材料复核。
\hypertarget{social:EQ-1902-CIXI-RETURN-BEIJING}{}\paragraph{光绪二十八年（1902年）}慈禧与光绪帝返回北京，战后朝廷开始推出新政措施。从地方治理、身份、家庭与物质生活的角度，这一节点要求追问它在何处改变了资源、道路、组织或日常风险。核心来源为《清德宗实录》光绪二十八年相关条；宫中朱批奏折、内阁大库题本与《清史稿》相关志传。它首先留下的是官府分类、任命、处置或资源调拨的痕迹；要判断基层是否照办，仍须寻找县档、账簿、契约或后续案件。账本所示的研究方向是宫廷政治、制度运作与中央—地方信息链研究。可核的条件包括辛丑议和结束流亡状态。应分层观察的后果是新政既回应危机，也服务朝廷重建合法性。证据界限：以光绪二十八年（1902年）的同期文书为定位起点；官修记录、奏报或外交文本服务特定机构立场，具体日期、规模、地方执行与对手视角须以独立材料复核。
\hypertarget{social:EQ-1902-NEW-POLICIES-EDUCATION}{}\paragraph{光绪二十八年（1902年）}清末新政在学制、军制、商政和地方行政方面逐步启动。从地方治理、身份、家庭与物质生活的角度，这一节点要求追问它在何处改变了资源、道路、组织或日常风险。核心来源为《清德宗实录》光绪二十八年相关条；内阁档、文集或报刊相关记录。这类编年材料可以标定朝廷获知和叙述事件的时间，却往往压低妇女、儿童、佃户、流民和非文书精英的行动。账本所示的研究方向是知识生产、教育、宗教与公共传播研究。可核的条件包括战败与财政外交压力显示旧制难以维持。应分层观察的后果是改革法令、试办与地方成效必须分别记录。证据界限：以光绪二十八年（1902年）的同期文书为定位起点；官修记录、奏报或外交文本服务特定机构立场，具体日期、规模、地方执行与对手视角须以独立材料复核。
\hypertarget{social:EQ-1903-BANNER-POSTS-REFORM}{}\paragraph{光绪二十九年（1903年）}清廷废除部分旗人官职垄断，调整旗籍制度与职务资格。从地方治理、身份、家庭与物质生活的角度，这一节点要求追问它在何处改变了资源、道路、组织或日常风险。核心来源为《清德宗实录》光绪二十九年相关条；地方志、督抚奏折及地方档案。这类编年材料可以标定朝廷获知和叙述事件的时间，却往往压低妇女、儿童、佃户、流民和非文书精英的行动。账本所示的研究方向是地方档案、迁徙、宗教、族群与社会动员研究。可核的条件包括八旗财政危机与官僚改革需求上升。应分层观察的后果是制度变化对不同旗籍人群的影响不均。证据界限：以光绪二十九年（1903年）的同期文书为定位起点；官修记录、奏报或外交文本服务特定机构立场，具体日期、规模、地方执行与对手视角须以独立材料复核。
\hypertarget{social:EQ-1903-SHANGHAI-COMMERCIAL-LAW}{}\paragraph{光绪二十九年（1903年）}商法、公司与工商管理的近代制度讨论在口岸和中央之间推进。从地方治理、身份、家庭与物质生活的角度，这一节点要求追问它在何处改变了资源、道路、组织或日常风险。核心来源为《清德宗实录》光绪二十九年相关条；户部奏销、盐政、河工或海关档案。它首先留下的是官府分类、任命、处置或资源调拨的痕迹；要判断基层是否照办，仍须寻找县档、账簿、契约或后续案件。账本所示的研究方向是财政账册、市场网络、灾害环境与区域经济研究。可核的条件包括国内商业、外贸和新式企业需要新的法律语言。应分层观察的后果是法典公布不等于商事审判和公司治理普遍落实。证据界限：以光绪二十九年（1903年）的同期文书为定位起点；官修记录、奏报或外交文本服务特定机构立场，具体日期、规模、地方执行与对手视角须以独立材料复核。
\hypertarget{social:EQ-1904-RUSSO-JAPANESE-WAR}{}\paragraph{光绪三十年（1904年）}日俄战争在中国东北爆发，清廷名义中立而东北遭受战场化。从地方治理、身份、家庭与物质生活的角度，这一节点要求追问它在何处改变了资源、道路、组织或日常风险。核心来源为《清德宗实录》光绪三十年相关条；军机处档、战事方略及地方奏报。这类编年材料可以标定朝廷获知和叙述事件的时间，却往往压低妇女、儿童、佃户、流民和非文书精英的行动。账本所示的研究方向是战争动员、军需、战场暴力与地方社会研究。可核的条件包括俄日争夺满洲铁路、港口和势力范围。应分层观察的后果是中国地方社会的损失与清廷的有限行动不能被列强叙事遮蔽。证据界限：以光绪三十年（1904年）的同期文书为定位起点；官修记录、奏报或外交文本服务特定机构立场，具体日期、规模、地方执行与对手视角须以独立材料复核。
\hypertarget{social:EQ-1904-IMPERIAL-EXAMINATION-DEBATE}{}\paragraph{光绪三十年（1904年）}科举存废与新式教育的争论加速。从地方治理、身份、家庭与物质生活的角度，这一节点要求追问它在何处改变了资源、道路、组织或日常风险。核心来源为《清德宗实录》光绪三十年相关条；内阁档、文集或报刊相关记录。这类编年材料可以标定朝廷获知和叙述事件的时间，却往往压低妇女、儿童、佃户、流民和非文书精英的行动。账本所示的研究方向是知识生产、教育、宗教与公共传播研究。可核的条件包括军政危机使人才选拔制度受到质疑。应分层观察的后果是讨论、诏令和实际学堂替代之间存在时间差。证据界限：以光绪三十年（1904年）的同期文书为定位起点；官修记录、奏报或外交文本服务特定机构立场，具体日期、规模、地方执行与对手视角须以独立材料复核。
\hypertarget{social:EQ-1905-ABOLITION-OF-CIVIL-SERVICE-EXAM}{}\paragraph{光绪三十一年（1905年）}清廷废除科举，延续千余年的取士制度终结。从地方治理、身份、家庭与物质生活的角度，这一节点要求追问它在何处改变了资源、道路、组织或日常风险。核心来源为《清德宗实录》光绪三十一年相关条；内阁档、文集或报刊相关记录。这类编年材料可以标定朝廷获知和叙述事件的时间，却往往压低妇女、儿童、佃户、流民和非文书精英的行动。账本所示的研究方向是知识生产、教育、宗教与公共传播研究。可核的条件包括新式教育、官僚改革与科举批评共同推动决策。应分层观察的后果是废科举不等于地方教育资源立即均衡或士人身份消失。证据界限：以光绪三十一年（1905年）的同期文书为定位起点；官修记录、奏报或外交文本服务特定机构立场，具体日期、规模、地方执行与对手视角须以独立材料复核。
\hypertarget{social:EQ-1905-TONGMENGHUI}{}\paragraph{光绪三十一年（1905年）}中国同盟会在东京成立，革命组织和宣传网络进一步整合。从地方治理、身份、家庭与物质生活的角度，这一节点要求追问它在何处改变了资源、道路、组织或日常风险。核心来源为《清德宗实录》光绪三十一年相关条；地方志、督抚奏折及地方档案。这类编年材料可以标定朝廷获知和叙述事件的时间，却往往压低妇女、儿童、佃户、流民和非文书精英的行动。账本所示的研究方向是地方档案、迁徙、宗教、族群与社会动员研究。可核的条件包括留日学生、海外华侨和反清团体形成跨境联系。应分层观察的后果是组织自称、成员规模与地方影响需要分别核验。证据界限：以光绪三十一年（1905年）的同期文书为定位起点；官修记录、奏报或外交文本服务特定机构立场，具体日期、规模、地方执行与对手视角须以独立材料复核。
\hypertarget{social:EQ-1906-CONSTITUTIONAL-PROMISE}{}\paragraph{光绪三十二年（1906年）}清廷宣布预备立宪，着手考察各国制度和改革官制。从地方治理、身份、家庭与物质生活的角度，这一节点要求追问它在何处改变了资源、道路、组织或日常风险。核心来源为《清德宗实录》光绪三十二年相关条；宫中朱批奏折、内阁大库题本与《清史稿》相关志传。它首先留下的是官府分类、任命、处置或资源调拨的痕迹；要判断基层是否照办，仍须寻找县档、账簿、契约或后续案件。账本所示的研究方向是宫廷政治、制度运作与中央—地方信息链研究。可核的条件包括新政需要重建统治合法性并回应立宪压力。应分层观察的后果是“预备”时间表和实际权力让渡之间存在重大距离。证据界限：以光绪三十二年（1906年）的同期文书为定位起点；官修记录、奏报或外交文本服务特定机构立场，具体日期、规模、地方执行与对手视角须以独立材料复核。
\hypertarget{social:EQ-1906-NEW-ARMY-REFORM}{}\paragraph{光绪三十二年（1906年）}新军编练与军事学堂扩展，军队成为新式政治动员的重要空间。从地方治理、身份、家庭与物质生活的角度，这一节点要求追问它在何处改变了资源、道路、组织或日常风险。核心来源为《清德宗实录》光绪三十二年相关条；军机处档、战事方略及地方奏报。这类编年材料可以标定朝廷获知和叙述事件的时间，却往往压低妇女、儿童、佃户、流民和非文书精英的行动。账本所示的研究方向是战争动员、军需、战场暴力与地方社会研究。可核的条件包括甲午、庚子和日俄战争显示旧军制弱点。应分层观察的后果是编制、装备和忠诚网络在各省差异很大。证据界限：以光绪三十二年（1906年）的同期文书为定位起点；官修记录、奏报或外交文本服务特定机构立场，具体日期、规模、地方执行与对手视角须以独立材料复核。
\hypertarget{social:EQ-1907-MANCHURIAN-PROVINCES}{}\paragraph{光绪三十三年（1907年）}清廷将东三省调整为行省体制，加强近代行政与边防经营。从地方治理、身份、家庭与物质生活的角度，这一节点要求追问它在何处改变了资源、道路、组织或日常风险。核心来源为《清德宗实录》光绪三十三年相关条；军机处档、理藩院档或相关方略。这类编年材料可以标定朝廷获知和叙述事件的时间，却往往压低妇女、儿童、佃户、流民和非文书精英的行动。账本所示的研究方向是多语边疆档案、盟旗/寺院/绿洲地方史与帝国治理研究。可核的条件包括日俄战争后东北主权和铁路问题尖锐化。应分层观察的后果是省制设立与铁路沿线列强控制并存。证据界限：以光绪三十三年（1907年）的同期文书为定位起点；官修记录、奏报或外交文本服务特定机构立场，具体日期、规模、地方执行与对手视角须以独立材料复核。
\hypertarget{social:EQ-1907-QING-LEGAL-REFORM}{}\paragraph{光绪三十三年（1907年）}清廷推进修律、审判和警政改革，尝试改造传统法律制度。从地方治理、身份、家庭与物质生活的角度，这一节点要求追问它在何处改变了资源、道路、组织或日常风险。核心来源为《清德宗实录》光绪三十三年相关条；宫中朱批奏折、内阁大库题本与《清史稿》相关志传。它首先留下的是官府分类、任命、处置或资源调拨的痕迹；要判断基层是否照办，仍须寻找县档、账簿、契约或后续案件。账本所示的研究方向是宫廷政治、制度运作与中央—地方信息链研究。可核的条件包括领事裁判权和国内治理压力推动法制改革。应分层观察的后果是新法文本、试行和基层司法实践不能混同。证据界限：以光绪三十三年（1907年）的同期文书为定位起点；官修记录、奏报或外交文本服务特定机构立场，具体日期、规模、地方执行与对手视角须以独立材料复核。
\hypertarget{social:EQ-1908-GUANGXU-CIXI-DEATHS}{}\paragraph{光绪三十四年（1908年）}光绪帝与慈禧太后相继去世，溥仪继位为宣统帝。从地方治理、身份、家庭与物质生活的角度，这一节点要求追问它在何处改变了资源、道路、组织或日常风险。核心来源为《清德宗实录》光绪三十四年相关条；宫中朱批奏折、内阁大库题本与《清史稿》相关志传。它首先留下的是官府分类、任命、处置或资源调拨的痕迹；要判断基层是否照办，仍须寻找县档、账簿、契约或后续案件。账本所示的研究方向是宫廷政治、制度运作与中央—地方信息链研究。可核的条件包括晚清权力核心突然更替。应分层观察的后果是摄政王载沣与幼帝体制面临改革和危机双重压力。证据界限：以光绪三十四年（1908年）的同期文书为定位起点；官修记录、奏报或外交文本服务特定机构立场，具体日期、规模、地方执行与对手视角须以独立材料复核。
\hypertarget{social:EQ-1908-CONSTITUTIONAL-OUTLINE}{}\paragraph{光绪三十四年（1908年）}清廷公布《钦定宪法大纲》等文件，明确预备立宪框架。从地方治理、身份、家庭与物质生活的角度，这一节点要求追问它在何处改变了资源、道路、组织或日常风险。核心来源为《清德宗实录》光绪三十四年相关条；宫中朱批奏折、内阁大库题本与《清史稿》相关志传。它首先留下的是官府分类、任命、处置或资源调拨的痕迹；要判断基层是否照办，仍须寻找县档、账簿、契约或后续案件。账本所示的研究方向是宫廷政治、制度运作与中央—地方信息链研究。可核的条件包括立宪压力与中央保留权力的意图并存。应分层观察的后果是文本承诺与咨议局、资政院实际权限存在落差。证据界限：以光绪三十四年（1908年）的同期文书为定位起点；官修记录、奏报或外交文本服务特定机构立场，具体日期、规模、地方执行与对手视角须以独立材料复核。
